%Supplementary material -- Springer Nature LaTeX template (sn-jnl)
\documentclass[pdflatex,sn-mathphys-num]{sn-jnl}

\usepackage{graphicx}%
\usepackage{amsmath,amssymb,amsfonts}%
\usepackage{xcolor}%
\usepackage{booktabs}%
\usepackage{longtable}%
\usepackage{tabularx}%
\usepackage{array}%

\newcommand{\RES}[1]{\textcolor{red}{[#1]}}
\newcolumntype{Y}{>{\centering\arraybackslash}X}
\renewcommand{\thetable}{S\arabic{table}}
\renewcommand{\thefigure}{S\arabic{figure}}

\raggedbottom

\begin{document}

\title[Supplementary Material]{Supplementary Material for: Hybrid text-anchored learning from heterogeneously annotated ECG databases for cross-database and open-vocabulary 12-lead ECG diagnosis}

\author*[1]{\fnm{Samiul} \sur{Alom}}\email{\RES{corresponding e-mail}}
\affil*[1]{\orgdiv{\RES{Department}}, \orgname{\RES{Institution}}, \orgaddress{\city{\RES{City}}, \country{\RES{Country}}}}

\abstract{This document contains the full diagnosis vocabulary with the frozen text descriptions used as label anchors (Table~S1), the database-by-diagnosis count matrix from which the annotated vocabularies are derived (Table~S2), the per-diagnosis results of the leave-one-database-out experiment (Table~S3), and the diagnoses removed in the annotation-removal experiment with their per-method results (Table~S4). All tables are generated by scripts from the preprocessing output and the logged run files.}

\maketitle

\section{Diagnosis vocabulary and text descriptions}\label{supp:descriptions}
Table~\ref{tab:descriptions} lists every diagnosis in the vocabulary with its SNOMED-CT code, whether it is one of the 26 classes scored in the PhysioNet/CinC 2021 Challenge, and the text description that is encoded once by the frozen BioLORD-2023 encoder to produce its anchor. Each anchor is formed as ``An ECG showing \textit{name}. \textit{Description}''; the name-only ablation uses ``An ECG showing \textit{name}.'' alone. Descriptions were written before any experiment and were not changed afterwards.

\begingroup\footnotesize
\begin{longtable}{@{}l l c p{0.62\textwidth}@{}}
\caption{Diagnosis vocabulary and frozen text descriptions used as label anchors. Scored = one of the 26 PhysioNet/CinC 2021 scored classes.}\label{tab:descriptions}\\
\toprule Abbr. & SNOMED-CT & Scored & Text anchor \\ \midrule \endfirsthead
\toprule Abbr. & SNOMED-CT & Scored & Text anchor \\ \midrule \endhead
AF & 164889003 & \checkmark & Atrial fibrillation: Irregularly irregular RR intervals with no distinct P waves; the baseline shows fibrillatory f waves and the ventricular rhythm varies beat to beat. \\
AFL & 164890007 & \checkmark & Atrial flutter: Regular atrial activity at about 250 to 350 per minute producing sawtooth flutter waves, best seen in leads II, III, aVF and V1, often with fixed 2:1 or variable atrioventricular conduction. \\
BBB & 6374002 & \checkmark & Bundle branch block: Delayed ventricular conduction through one bundle branch giving a widened QRS complex of 120 ms or more with abnormal terminal QRS morphology. \\
Brady & 426627000 & \checkmark & Bradycardia: Ventricular rate below 60 beats per minute regardless of the underlying rhythm. \\
LBBB & 164909002 & \checkmark & Left bundle branch block: QRS duration of 120 ms or more with broad notched R waves in I, aVL, V5 and V6, absent septal q waves in lateral leads and deep S waves in V1, with discordant ST-T changes. \\
RBBB & 59118001 & \checkmark & Right bundle branch block: QRS duration of 120 ms or more with an rsR' pattern in V1 and V2 and a wide slurred S wave in leads I and V6. \\
IAVB & 270492004 & \checkmark & First degree atrioventricular block: Every P wave is conducted to the ventricles but the PR interval is prolonged beyond 200 ms. \\
IRBBB & 713426002 & \checkmark & Incomplete right bundle branch block: rsR' pattern in V1 and V2 resembling right bundle branch block but with QRS duration between 110 and 120 ms. \\
LAD & 39732003 & \checkmark & Left axis deviation: Mean frontal plane QRS axis more negative than minus 30 degrees: positive QRS in lead I and predominantly negative QRS in leads II and aVF. \\
LAnFB & 445118002 & \checkmark & Left anterior fascicular block: Left axis deviation between minus 45 and minus 90 degrees with qR in I and aVL, rS in II, III and aVF, and QRS duration under 120 ms. \\
LQRSV & 251146004 & \checkmark & Low QRS voltages: Small QRS amplitudes: peak-to-peak below 0.5 mV in all limb leads or below 1.0 mV in all precordial leads. \\
NSIVCB & 698252002 & \checkmark & Nonspecific intraventricular conduction disturbance: QRS duration above 110 ms without the morphology criteria of left or right bundle branch block. \\
NSR & 426783006 & \checkmark & Sinus rhythm: Normal sinus rhythm: upright P wave in leads I and II before every QRS complex, constant PR interval and a regular rate between 60 and 100 beats per minute. \\
SVPB & 63593006 & \checkmark & Supraventricular premature beats: Premature atrial or junctional beats: early narrow QRS complexes preceded by an abnormal or hidden P wave, usually followed by a non-compensatory pause. \\
PR & 10370003 & \checkmark & Pacing rhythm: Electronic pacemaker rhythm with sharp pacing spikes preceding paced P waves or wide paced QRS complexes. \\
PRWP & 365413008 & \checkmark & Poor R wave progression: R wave amplitude fails to increase normally across the precordial leads, remaining 3 mm or less in V3, with delayed transition zone. \\
VPB & 17338001 & \checkmark & Ventricular premature beats: Premature wide and bizarre QRS complexes not preceded by a P wave, with discordant T waves and usually a full compensatory pause. \\
LPR & 164947007 & \checkmark & Prolonged PR interval: PR interval longer than 200 ms, measured from the onset of the P wave to the onset of the QRS complex. \\
LQT & 111975006 & \checkmark & Prolonged QT interval: Heart-rate corrected QT interval prolonged, above about 450 ms in men or 460 ms in women, reflecting delayed ventricular repolarization. \\
QAb & 164917005 & \checkmark & Q wave abnormal: Pathological Q waves wider than 40 ms or deeper than one quarter of the following R wave, suggesting prior myocardial necrosis. \\
RAD & 47665007 & \checkmark & Right axis deviation: Mean frontal plane QRS axis greater than plus 90 degrees: negative QRS in lead I and positive QRS in lead aVF. \\
SA & 427393009 & \checkmark & Sinus arrhythmia: Sinus P waves with normal morphology but phasic variation of the PP interval by more than 0.12 seconds, typically with respiration. \\
SB & 426177001 & \checkmark & Sinus bradycardia: Sinus rhythm with normal P waves before each QRS complex at a rate below 60 beats per minute. \\
STach & 427084000 & \checkmark & Sinus tachycardia: Sinus rhythm with normal P waves before each QRS complex at a rate above 100 beats per minute. \\
TAb & 164934002 & \checkmark & T wave abnormal: Abnormal T waves that are flattened, low amplitude, biphasic or inverted in leads where they are normally upright, indicating abnormal repolarization. \\
TInv & 59931005 & \checkmark & T wave inversion: Negative T waves in leads where the T wave is normally positive, such as I, II and V3 to V6. \\
MI & 164865005 &  & Myocardial infarction: Evidence of myocardial infarction: pathological Q waves, loss of R wave amplitude and ST-T changes in contiguous leads matching a coronary territory. \\
LVH & 164873001 &  & Left ventricular hypertrophy: Increased QRS voltage over the left ventricle, for example S in V1 plus R in V5 or V6 above 35 mm, often with lateral ST depression and T inversion of strain pattern. \\
STD & 429622005 &  & ST depression: Horizontal or downsloping depression of the ST segment of at least 0.5 mm below the PR baseline in two contiguous leads. \\
STE & 164931005 &  & ST elevation: Elevation of the ST segment at the J point above the baseline in two or more contiguous leads. \\
LVHV & 55827005 &  & Left ventricular high voltage: Increased QRS amplitudes in left precordial or limb leads meeting voltage criteria without other features of hypertrophy. \\
IIAVB & 195042002 &  & Second degree atrioventricular block: Intermittent failure of atrioventricular conduction: some P waves are not followed by a QRS complex, with or without progressive PR prolongation. \\
AVB & 233917008 &  & Atrioventricular block: Impaired conduction from atria to ventricles shown by prolonged PR interval or P waves not followed by QRS complexes. \\
NSSTTA & 428750005 &  & Nonspecific ST-T abnormality: Minor ST segment deviation or T wave flattening that does not meet criteria for ischemia, injury or a specific diagnosis. \\
LAE & 67741000119109 &  & Left atrial enlargement: Broad notched P wave longer than 120 ms in lead II with a deep terminal negative P component in V1. \\
RAH & 446358003 &  & Right atrial hypertrophy: Tall peaked P waves above 2.5 mm in the inferior leads indicating right atrial enlargement. \\
RVH & 89792004 &  & Right ventricular hypertrophy: Right axis deviation with dominant R wave in V1 and deep S waves in V5 and V6 indicating right ventricular hypertrophy. \\
ATach & 713422000 &  & Atrial tachycardia: Regular narrow complex tachycardia above 100 per minute with non-sinus P waves of uniform abnormal morphology. \\
SVT & 426761007 &  & Supraventricular tachycardia: Rapid regular narrow QRS tachycardia above 150 beats per minute originating above the ventricles, with P waves hidden or retrograde. \\
CCR & 251198002 &  & Clockwise cardiac rotation: Delayed precordial transition with the R/S ratio becoming greater than one only in V5 or V6. \\
CCWR & 251199005 &  & Counterclockwise cardiac rotation: Early precordial transition with dominant R waves already in V1 or V2. \\
PWC & 164912004 &  & P wave change: Abnormal P wave morphology, amplitude or axis compared with normal sinus P waves. \\
FQRS & 164942001 &  & Fragmented QRS: Additional notches or RSR' patterns within a narrow QRS complex in two contiguous leads. \\
STIAb & 164930006 &  & ST interval abnormal: Abnormal ST segment shape or level not meeting specific elevation or depression criteria. \\
UAb & 164937009 &  & U wave abnormal: Prominent or inverted U waves following the T wave. \\
\bottomrule \end{longtable}
\endgroup


\section{Per-database label counts}\label{supp:counts}
Table~\ref{tab:counts} gives the number of positive recordings of each vocabulary diagnosis in each database after the official code equivalences are applied. A database annotates a diagnosis when the count is at least ten (shown in bold); the vocabulary contains the diagnoses with at least 50 positives in total and at least ten in one database.

{\footnotesize\begin{longtable}{@{}llrrrrr@{}}
\caption{Positive recordings per diagnosis and database after the official code equivalences, i.e.\ the label matrix used for training and evaluation. Bold: the database annotates the diagnosis ($\geq10$ positives).}\label{tab:counts}\\
\toprule Abbr. & SNOMED-CT & CPSC & PTB-XL & Georgia & S--N & Total \\ \midrule \endfirsthead
\toprule Abbr. & SNOMED-CT & CPSC & PTB-XL & Georgia & S--N & Total \\ \midrule \endhead
\bottomrule \endfoot
NSR & 426783006 & \textbf{922} & \textbf{18{,}092} & \textbf{1{,}752} & \textbf{8{,}125} & 28{,}891 \\
SB & 426177001 & \textbf{45} & \textbf{637} & \textbf{1{,}677} & \textbf{16{,}559} & 18{,}918 \\
TAb & 164934002 & \textbf{22} & \textbf{2{,}345} & \textbf{2{,}306} & \textbf{7{,}043} & 11{,}716 \\
STach & 427084000 & \textbf{303} & \textbf{826} & \textbf{1{,}261} & \textbf{7{,}255} & 9{,}645 \\
AFL & 164890007 & \textbf{54} & \textbf{73} & \textbf{186} & \textbf{8{,}060} & 8{,}373 \\
LAD & 39732003 & 0 & \textbf{5{,}146} & \textbf{940} & \textbf{1{,}545} & 7{,}631 \\
MI & 164865005 & \textbf{374} & \textbf{5{,}261} & 7 & \textbf{123} & 5{,}765 \\
LVHV & 55827005 & 0 & 0 & 0 & \textbf{5{,}401} & 5{,}401 \\
AF & 164889003 & \textbf{1{,}374} & \textbf{1{,}514} & \textbf{570} & \textbf{1{,}780} & 5{,}238 \\
RBBB & 59118001 & \textbf{1{,}971} & \textbf{542} & \textbf{556} & \textbf{1{,}745} & 4{,}814 \\
NSSTTA & 428750005 & \textbf{1{,}290} & \textbf{381} & \textbf{1{,}883} & \textbf{1{,}158} & 4{,}712 \\
LVH & 164873001 & \textbf{158} & \textbf{2{,}359} & \textbf{1{,}232} & \textbf{647} & 4{,}396 \\
TInv & 59931005 & 5 & \textbf{294} & \textbf{812} & \textbf{2{,}877} & 3{,}988 \\
SA & 427393009 & \textbf{11} & \textbf{772} & \textbf{455} & \textbf{2{,}550} & 3{,}788 \\
STD & 429622005 & \textbf{926} & \textbf{1{,}009} & \textbf{38} & \textbf{1{,}668} & 3{,}641 \\
IAVB & 270492004 & \textbf{828} & \textbf{797} & \textbf{769} & \textbf{1{,}140} & 3{,}534 \\
SVPB & 63593006 & \textbf{740} & \textbf{555} & \textbf{640} & \textbf{1{,}321} & 3{,}256 \\
STIAb & 164930006 & \textbf{481} & 0 & \textbf{992} & \textbf{801} & 2{,}274 \\
LAnFB & 445118002 & 0 & \textbf{1{,}626} & \textbf{180} & \textbf{380} & 2{,}186 \\
LQRSV & 251146004 & 0 & \textbf{182} & \textbf{374} & \textbf{1{,}613} & 2{,}169 \\
QAb & 164917005 & 1 & \textbf{548} & \textbf{464} & \textbf{1{,}063} & 2{,}076 \\
VPB & 17338001 & \textbf{194} & 0 & \textbf{357} & \textbf{1{,}385} & 1{,}936 \\
LQT & 111975006 & 4 & \textbf{118} & \textbf{1{,}391} & \textbf{391} & 1{,}904 \\
IRBBB & 713426002 & \textbf{86} & \textbf{1{,}118} & \textbf{407} & \textbf{246} & 1{,}857 \\
NSIVCB & 698252002 & 4 & \textbf{789} & \textbf{203} & \textbf{771} & 1{,}767 \\
PR & 10370003 & 3 & \textbf{296} & 0 & \textbf{1{,}182} & 1{,}481 \\
LAE & 67741000119109 & 1 & \textbf{427} & \textbf{870} & 1 & 1{,}299 \\
LBBB & 164909002 & \textbf{274} & \textbf{536} & \textbf{231} & \textbf{240} & 1{,}281 \\
RAD & 47665007 & 1 & \textbf{343} & \textbf{83} & \textbf{853} & 1{,}280 \\
SVT & 426761007 & 3 & \textbf{27} & \textbf{32} & \textbf{724} & 786 \\
STE & 164931005 & \textbf{286} & \textbf{28} & \textbf{134} & \textbf{176} & 624 \\
BBB & 6374002 & 0 & 0 & \textbf{116} & \textbf{385} & 501 \\
LPR & 164947007 & 0 & \textbf{340} & 0 & \textbf{52} & 392 \\
RVH & 89792004 & \textbf{20} & \textbf{126} & \textbf{86} & \textbf{110} & 342 \\
ATach & 713422000 & \textbf{15} & 0 & \textbf{28} & \textbf{297} & 340 \\
AVB & 233917008 & 5 & 0 & \textbf{74} & \textbf{244} & 323 \\
Brady & 426627000 & \textbf{271} & 0 & 6 & 7 & 284 \\
CCWR & 251199005 & 0 & 0 & 0 & \textbf{162} & 162 \\
RAH & 446358003 & \textbf{18} & \textbf{99} & 0 & \textbf{36} & 153 \\
PWC & 164912004 & 0 & 0 & 0 & \textbf{142} & 142 \\
UAb & 164937009 & 1 & 0 & 0 & \textbf{136} & 137 \\
IIAVB & 195042002 & \textbf{21} & \textbf{14} & \textbf{23} & \textbf{66} & 124 \\
CCR & 251198002 & 0 & 0 & 0 & \textbf{76} & 76 \\
\end{longtable}
}

\section{Per-diagnosis leave-one-database-out results}\label{supp:perlabel}
Table~\ref{tab:perlabel} reports, for each held-out database and each evaluable diagnosis, the number of positives, the AUROC of the proposed method and of the strongest baseline of that database (the baseline with the highest seed-0 macro AUROC on the same labels), and the Holm-adjusted DeLong $p$-value of their difference. Unseen diagnoses are appended per database.

{\footnotesize\begin{longtable}{@{}llrrrrl@{}}
\caption{Per-diagnosis AUROC ($\times100$) under leave-one-database-out (seed 0): proposed method against the strongest baseline of each held-out database, with DeLong $p$ Holm-adjusted over all per-diagnosis tests. Unseen diagnoses ($\dagger$) are listed with the best baseline by AUROC and no test.}\label{tab:perlabel}\\
\toprule Target & Diagnosis & Positives & Ours & Baseline & $p_{\mathrm{Holm}}$ & Baseline \\ \midrule \endfirsthead
\toprule Target & Diagnosis & Positives & Ours & Baseline & $p_{\mathrm{Holm}}$ & Baseline \\ \midrule \endhead
\bottomrule \endfoot
CPSC & AF & 1374 & 96.6 & 97.3 & 0.292 & GroupDRO \\
CPSC & AFL & 54 & 94.8 & 96.0 & 1.000 & GroupDRO \\
CPSC & LBBB & 274 & 95.9 & 95.3 & 1.000 & GroupDRO \\
CPSC & RBBB & 1971 & 95.4 & 94.4 & \textbf{$<$0.001} & GroupDRO \\
CPSC & IAVB & 828 & 96.8 & 96.3 & 1.000 & GroupDRO \\
CPSC & IRBBB & 86 & 75.7 & 76.5 & 1.000 & GroupDRO \\
CPSC & NSR & 922 & 84.9 & 88.9 & \textbf{$<$0.001} & GroupDRO \\
CPSC & SVPB & 740 & 91.4 & 91.8 & 1.000 & GroupDRO \\
CPSC & VPB & 194 & 83.8 & 84.7 & 1.000 & GroupDRO \\
CPSC & SA & 11 & 80.0 & 78.3 & 1.000 & GroupDRO \\
CPSC & SB & 45 & 89.3 & 89.9 & 1.000 & GroupDRO \\
CPSC & STach & 303 & 96.2 & 96.3 & 1.000 & GroupDRO \\
CPSC & TAb & 22 & 79.0 & 78.3 & 1.000 & GroupDRO \\
CPSC & MI & 374 & 81.7 & 81.3 & 1.000 & GroupDRO \\
CPSC & LVH & 158 & 90.2 & 92.2 & 0.076 & GroupDRO \\
CPSC & STD & 926 & 65.9 & 61.7 & \textbf{$<$0.001} & GroupDRO \\
CPSC & STE & 286 & 81.1 & 82.5 & 1.000 & GroupDRO \\
CPSC & IIAVB & 21 & 94.9 & 92.2 & 1.000 & GroupDRO \\
CPSC & NSSTTA & 1290 & 75.8 & 73.5 & \textbf{$<$0.001} & GroupDRO \\
CPSC & RAH & 18 & 87.2 & 87.3 & 1.000 & GroupDRO \\
CPSC & RVH & 20 & 84.3 & 82.1 & 1.000 & GroupDRO \\
CPSC & ATach & 15 & 76.4 & 83.3 & 1.000 & GroupDRO \\
CPSC & STIAb & 481 & 44.6 & 47.3 & 0.262 & GroupDRO \\
CPSC & Brady$^\dagger$ & 271 & 94.6 & 89.9 & -- & GroupDRO \\
\midrule
PTB-XL & AF & 1514 & 98.7 & 98.7 & 1.000 & Li-style \\
PTB-XL & AFL & 73 & 97.2 & 96.9 & 1.000 & Li-style \\
PTB-XL & LBBB & 536 & 99.5 & 99.5 & 1.000 & Li-style \\
PTB-XL & RBBB & 542 & 99.3 & 99.3 & 0.500 & Li-style \\
PTB-XL & IAVB & 797 & 95.2 & 95.4 & 1.000 & Li-style \\
PTB-XL & IRBBB & 1118 & 93.2 & 92.9 & 1.000 & Li-style \\
PTB-XL & LAD & 5146 & 86.5 & 87.0 & \textbf{$<$0.001} & Li-style \\
PTB-XL & LAnFB & 1626 & 97.6 & 97.5 & 1.000 & Li-style \\
PTB-XL & LQRSV & 182 & 79.4 & 79.1 & 1.000 & Li-style \\
PTB-XL & NSIVCB & 789 & 75.4 & 74.9 & 1.000 & Li-style \\
PTB-XL & NSR & 18092 & 85.7 & 86.1 & \textbf{$<$0.001} & Li-style \\
PTB-XL & SVPB & 555 & 88.2 & 87.6 & 1.000 & Li-style \\
PTB-XL & PR & 296 & 96.5 & 96.4 & 1.000 & Li-style \\
PTB-XL & LPR & 340 & 95.2 & 95.4 & 1.000 & Li-style \\
PTB-XL & LQT & 118 & 91.5 & 91.5 & 1.000 & Li-style \\
PTB-XL & QAb & 548 & 72.6 & 71.5 & 1.000 & Li-style \\
PTB-XL & RAD & 343 & 96.6 & 96.7 & 1.000 & Li-style \\
PTB-XL & SA & 772 & 92.3 & 92.6 & 1.000 & Li-style \\
PTB-XL & SB & 637 & 93.5 & 93.5 & 1.000 & Li-style \\
PTB-XL & STach & 826 & 99.0 & 99.0 & 1.000 & Li-style \\
PTB-XL & TAb & 2345 & 71.1 & 71.7 & \textbf{0.037} & Li-style \\
PTB-XL & TInv & 294 & 89.6 & 89.1 & 1.000 & Li-style \\
PTB-XL & MI & 5261 & 82.9 & 81.8 & \textbf{$<$0.001} & Li-style \\
PTB-XL & LVH & 2359 & 87.1 & 87.0 & 1.000 & Li-style \\
PTB-XL & STD & 1009 & 78.0 & 79.0 & \textbf{0.033} & Li-style \\
PTB-XL & STE & 28 & 51.0 & 49.5 & 1.000 & Li-style \\
PTB-XL & IIAVB & 14 & 94.3 & 89.2 & 1.000 & Li-style \\
PTB-XL & NSSTTA & 381 & 75.9 & 76.1 & 1.000 & Li-style \\
PTB-XL & LAE & 427 & 80.2 & 79.3 & 1.000 & Li-style \\
PTB-XL & RAH & 99 & 79.5 & 82.8 & 1.000 & Li-style \\
PTB-XL & RVH & 126 & 87.1 & 84.2 & 0.320 & Li-style \\
PTB-XL & SVT & 27 & 98.6 & 98.1 & 1.000 & Li-style \\
\midrule
Georgia & AF & 570 & 92.9 & 93.0 & 1.000 & Masked \\
Georgia & AFL & 186 & 91.2 & 92.0 & 1.000 & Masked \\
Georgia & BBB & 116 & 90.3 & 78.1 & \textbf{$<$0.001} & Masked \\
Georgia & LBBB & 231 & 97.7 & 98.0 & 1.000 & Masked \\
Georgia & RBBB & 556 & 96.9 & 96.6 & 1.000 & Masked \\
Georgia & IAVB & 769 & 97.2 & 97.2 & 1.000 & Masked \\
Georgia & IRBBB & 407 & 91.9 & 91.6 & 1.000 & Masked \\
Georgia & LAD & 940 & 94.8 & 94.6 & 1.000 & Masked \\
Georgia & LAnFB & 180 & 96.1 & 96.3 & 1.000 & Masked \\
Georgia & LQRSV & 374 & 85.4 & 87.1 & 0.075 & Masked \\
Georgia & NSIVCB & 203 & 81.5 & 80.3 & 1.000 & Masked \\
Georgia & NSR & 1752 & 87.3 & 87.1 & 1.000 & Masked \\
Georgia & SVPB & 640 & 90.4 & 91.2 & 1.000 & Masked \\
Georgia & VPB & 357 & 85.6 & 85.8 & 1.000 & Masked \\
Georgia & LQT & 1391 & 77.8 & 81.3 & \textbf{$<$0.001} & Masked \\
Georgia & QAb & 464 & 51.9 & 52.6 & 1.000 & Masked \\
Georgia & RAD & 83 & 96.8 & 96.2 & 1.000 & Masked \\
Georgia & SA & 455 & 96.7 & 96.8 & 1.000 & Masked \\
Georgia & SB & 1677 & 99.2 & 99.1 & 1.000 & Masked \\
Georgia & STach & 1261 & 99.3 & 99.1 & 1.000 & Masked \\
Georgia & TAb & 2306 & 78.5 & 75.4 & \textbf{$<$0.001} & Masked \\
Georgia & TInv & 812 & 68.0 & 68.3 & 1.000 & Masked \\
Georgia & LVH & 1232 & 83.9 & 83.3 & 1.000 & Masked \\
Georgia & STD & 38 & 57.9 & 60.9 & 1.000 & Masked \\
Georgia & STE & 134 & 80.8 & 80.8 & 1.000 & Masked \\
Georgia & IIAVB & 23 & 93.5 & 94.0 & 1.000 & Masked \\
Georgia & AVB & 74 & 92.5 & 96.8 & \textbf{$<$0.001} & Masked \\
Georgia & NSSTTA & 1883 & 64.9 & 65.5 & 0.767 & Masked \\
Georgia & LAE & 870 & 67.5 & 68.8 & 0.706 & Masked \\
Georgia & RVH & 86 & 90.2 & 91.3 & 1.000 & Masked \\
Georgia & ATach & 28 & 94.5 & 94.8 & 1.000 & Masked \\
Georgia & SVT & 32 & 94.8 & 97.0 & 1.000 & Masked \\
Georgia & STIAb & 992 & 60.5 & 58.1 & \textbf{0.002} & Masked \\
\midrule
S--N & AF & 1780 & 92.9 & 92.7 & 1.000 & Masked \\
S--N & AFL & 8060 & 83.6 & 77.9 & \textbf{$<$0.001} & Masked \\
S--N & BBB & 385 & 77.9 & 85.8 & \textbf{$<$0.001} & Masked \\
S--N & LBBB & 240 & 93.9 & 93.1 & 1.000 & Masked \\
S--N & RBBB & 1745 & 99.2 & 99.2 & 1.000 & Masked \\
S--N & IAVB & 1140 & 97.0 & 96.2 & \textbf{$<$0.001} & Masked \\
S--N & IRBBB & 246 & 94.5 & 94.6 & 1.000 & Masked \\
S--N & LAD & 1545 & 97.3 & 97.1 & \textbf{$<$0.001} & Masked \\
S--N & LAnFB & 380 & 97.8 & 97.7 & 1.000 & Masked \\
S--N & LQRSV & 1613 & 69.7 & 72.5 & \textbf{$<$0.001} & Masked \\
S--N & NSIVCB & 771 & 67.8 & 61.2 & \textbf{$<$0.001} & Masked \\
S--N & NSR & 8125 & 87.4 & 88.2 & \textbf{$<$0.001} & Masked \\
S--N & SVPB & 1321 & 95.1 & 94.3 & 0.262 & Masked \\
S--N & PR & 1182 & 93.1 & 92.3 & 1.000 & Masked \\
S--N & VPB & 1385 & 87.8 & 87.1 & 0.706 & Masked \\
S--N & LPR & 52 & 92.2 & 91.5 & 1.000 & Masked \\
S--N & LQT & 391 & 88.8 & 88.8 & 1.000 & Masked \\
S--N & QAb & 1063 & 79.0 & 78.3 & 1.000 & Masked \\
S--N & RAD & 853 & 95.4 & 95.5 & 1.000 & Masked \\
S--N & SA & 2550 & 92.3 & 93.0 & \textbf{0.021} & Masked \\
S--N & SB & 16559 & 99.0 & 98.9 & \textbf{$<$0.001} & Masked \\
S--N & STach & 7255 & 98.5 & 98.6 & 0.324 & Masked \\
S--N & TAb & 7043 & 78.1 & 80.8 & \textbf{$<$0.001} & Masked \\
S--N & TInv & 2877 & 88.1 & 87.1 & \textbf{0.002} & Masked \\
S--N & MI & 123 & 92.4 & 90.7 & 1.000 & Masked \\
S--N & LVH & 647 & 91.3 & 90.5 & 0.110 & Masked \\
S--N & STD & 1668 & 72.2 & 72.1 & 1.000 & Masked \\
S--N & STE & 176 & 76.4 & 77.9 & 1.000 & Masked \\
S--N & IIAVB & 66 & 92.3 & 91.6 & 1.000 & Masked \\
S--N & AVB & 244 & 90.5 & 91.3 & 1.000 & Masked \\
S--N & NSSTTA & 1158 & 76.4 & 76.6 & 1.000 & Masked \\
S--N & RAH & 36 & 93.0 & 94.8 & 1.000 & Masked \\
S--N & RVH & 110 & 95.4 & 93.6 & 1.000 & Masked \\
S--N & ATach & 297 & 85.6 & 83.8 & \textbf{0.028} & Masked \\
S--N & SVT & 724 & 98.7 & 98.7 & 1.000 & Masked \\
S--N & STIAb & 801 & 59.4 & 58.8 & 1.000 & Masked \\
S--N & LVHV$^\dagger$ & 5401 & 71.4 & 47.6 & -- & Naive \\
S--N & CCR$^\dagger$ & 76 & 79.2 & 81.5 & -- & CORAL \\
S--N & CCWR$^\dagger$ & 162 & 55.1 & 61.5 & -- & Naive \\
S--N & PWC$^\dagger$ & 142 & 59.6 & 71.7 & -- & Naive \\
S--N & UAb$^\dagger$ & 136 & 76.7 & 44.5 & -- & DANN \\

\end{longtable}
}

\section{Diagnoses removed in the annotation-removal experiment}\label{supp:e2}
Table~\ref{tab:e2diag} lists the PTB-XL diagnoses removed from training supervision at each removal rate, drawn once with a fixed seed and shared by all methods, with the test AUROC of each method on each removed diagnosis.

\begin{table}[h]\centering\footnotesize
\caption{Diagnoses removed from PTB-XL training supervision in the annotation-removal experiment and their test AUROC ($\times100$, single seed). Positives: number of positive test recordings.}\label{tab:e2diag}
\begin{tabular}{llrrrr}
\toprule $r$ & Diagnosis & Positives & Naive & Li-style & Ours \\ \midrule
0.25 & AFL & 10 & 51.2 & 79.8 & 91.8 \\
0.25 & LAnFB & 163 & 54.0 & 63.4 & 92.6 \\
0.25 & SVPB & 65 & 73.8 & 44.4 & 79.8 \\
0.25 & PR & 26 & 26.7 & 93.4 & 69.8 \\
0.25 & RAD & 36 & 37.3 & 61.4 & 47.8 \\
0.25 & STach & 75 & 40.8 & 12.0 & 15.5 \\
0.25 & TInv & 30 & 46.2 & 36.6 & 69.1 \\
\midrule
0.5 & AF & 149 & 49.9 & 49.4 & 77.2 \\
0.5 & AFL & 10 & 67.5 & 44.2 & 87.1 \\
0.5 & LBBB & 55 & 33.3 & 30.6 & 52.4 \\
0.5 & RBBB & 49 & 65.7 & 53.4 & 93.5 \\
0.5 & IAVB & 77 & 66.3 & 29.5 & 78.2 \\
0.5 & LAnFB & 163 & 62.7 & 73.7 & 94.4 \\
0.5 & LQRSV & 20 & 59.2 & 68.4 & 49.6 \\
0.5 & SVPB & 65 & 66.3 & 39.1 & 73.3 \\
0.5 & PR & 26 & 57.9 & 90.5 & 67.8 \\
0.5 & LPR & 37 & 57.0 & 35.6 & 73.8 \\
0.5 & LQT & 15 & 57.9 & 46.3 & 64.4 \\
0.5 & RAD & 36 & 61.2 & 66.9 & 37.7 \\
0.5 & SA & 77 & 50.7 & 34.1 & 52.5 \\
0.5 & STach & 75 & 24.4 & 16.8 & 43.6 \\
0.5 & TAb & 229 & 62.0 & 52.7 & 63.8 \\
0.5 & TInv & 30 & 42.4 & 23.3 & 82.5 \\
0.5 & NSSTTA & 42 & 59.9 & 54.9 & 82.1 \\
\bottomrule
\end{tabular}
\end{table}


\end{document}
